\begin{sloppypar}
, Supercritical Helium Loop, and Dual-Stage TEG Power GenerationThermal management of the reaction chamber first-wall armor tiles and high-heat-flux divertors is provided by a closed~$10.0$\,MPa supercritical helium primary cooling circuit. The total primary coolant mass flow rate is $\dot{m}_{\text{total}} = 450.0\text{ kg/s}$, entering the chamber distributor plenums at $T_{\text{in}} = 300.0\text{ K}$ and returning from the heat-exchanger manifold at $T_{\text{out}} = 900.0\text{ K}$. Operating at an average isobaric heat capacity $c_p = 5,193.15\text{ J/kg}\cdot\text{K}$, the nominal steady-state thermal extraction capacity is $\dot{Q}_{\text{th}} = 1,402.15\text{ MW}_{\text{th}}$.Coolant routing distributes the primary flow across 64 parallel micro-channel heat exchanger sectors, each comprising smooth channels with hydraulic diameter $D_h = 5.35\text{ mm}$ and length $L = 4.2\text{ m}$. Compressible Reynolds-Averaged Navier-Stokes formulations utilizing the $k-\omega$ SST turbulence model resolve steep property variations across the fluid domain. Under nominal operating conditions, the integrated total loop pressure drop is:\[
\resizebox{\columnwidth}{!}{$\displaystyle \Delta P_{\text{total}} = \Delta P_{\text{friction}} + \Delta P_{\text{acceleration}} + \Delta P_{\text{form}} = 283.4\text{ kPa}$}
\]
Circulation is driven by a multi-stage centrifugal compressor operating at inlet conditions ($300\text{ K}$, $10.0\text{ MPa}$, $\rho_{\text{in}} = 16.05\text{ kg/m}^3$) with a net electrical-to-fluid efficiency $\eta_{\text{pump}} = 0.8624$. The required electrical pumping power is:\[
\resizebox{\columnwidth}{!}{$\displaystyle W_{\text{pump}} = \frac{\dot{m} \Delta P_{\text{total}}}{\rho_{\text{in}} \eta_{\text{pump}}} = 9.213\text{ MW}_{\text{elec}}$}
\]
This satisfies the operating threshold $W_{\text{pump}} \le 15.0\text{ MW}_{\text{elec}}$ with a $38.6\%$ reserve margin.During each 100 Hz ignition event, the $10.0\,\mu\text{s}$ surface flash injects an intense thermal pulse into the channel walls. Acoustic speed in the supercritical fluid ranges from $1,028\text{ m/s}$ (inlet) to $1,762\text{ m/s}$ (outlet). The resulting Joukowsky-thermal acoustic pressure wave amplitude is bounded by $\Delta p_{\text{surge}} = 1.44\text{ kPa}$, avoiding high-frequency hydraulic shock and preventing mechanical channel wall delamination.Thermal energy extracted by the supercritical helium loop is routed to a primary solid-state energy conversion system consisting of a dual-stage thermoelectric generator network. The upper stage ($600\text{ K} \to 900\text{ K}$) utilizes Half-Heusler alloy modules ($\text{HfZrNiSnSb}$ / $\text{FeNbTiSb}$), while the lower stage ($300\text{ K} \to 600\text{ K}$) employs Skutterudite modules ($\text{BaLaYbCo}_4\text{Sb}_{12}$ / $\text{CeFeCoSb}_{12}$). Across the 600 K temperature differential, the cascaded network achieves a net conversion efficiency $\eta_{\text{TEG}} = 33.804\%$ while processing a thermal input $\dot{Q}_{\text{TEG, th}} = 1,325.000\text{ MW}_{\text{th}}$. The system produces a continuous electrical output $P_{\text{elec}} = 447.903\text{ MW}_{\text{elec}}$. The TEG array is directly coupled to the 1,800-module LANR starter grid from shbt-cf, providing $358.32\text{ kA}$ at $1.25\text{ kV}$ into an equivalent matched bus impedance $R_{\text{load}} = 3.488\text{ m}\Omega$.\end{sloppypar}
